\section{Space Complexity --- Machines --- Bank}
%%% generated by the blueprint pipeline from TCSlib.Complexity.SpaceComplexity.Machines.Bank
\subsection*{Space Complexity --- Machines --- Bank: Introduction}

Declarations of the Lean module \texttt{TCSlib.Complexity.SpaceComplexity.Machines.Bank}.

\subsection*{Space Complexity --- Machines --- Bank: Results}

\begin{definition}[Work tape padding]
\lean{Complexity.LogProg.padTM}
\label{Complexity.LogProg.padTM}
\leanok
\uses{Turing.MultiTapeTM}
Given a deterministic machine with $k$ Boolean work tapes and an integer $K$, define a
machine with $K$ work tapes and the same initial state and state set. Its transition
simulates the original machine using the first $k$ tapes when available, treats any
missing original tape as blank, and leaves every extra tape blank and stationary.
\end{definition}

\begin{definition}[Configuration padding]
\lean{Complexity.LogProg.padCfg}
\label{Complexity.LogProg.padCfg}
\leanok
\uses{Turing.Cfg}
Given a configuration on $k$ Boolean work tapes and a natural number $K$, define a
configuration on $K$ work tapes with the same state, input head position, and output.
Its first $k$ work tapes and head positions are retained, while any additional tapes are
blank and have their heads at position zero.
\end{definition}

\begin{lemma}[Padding preserves one machine step]
\lean{Complexity.LogProg.padTM_step}
\label{Complexity.LogProg.padTM_step}
\leanok
\uses{Complexity.LogProg.padCfg, Complexity.LogProg.padTM, Turing.Action.apply, Turing.Cfg, Turing.Cfg.inputSymbol, Turing.Cfg.workTapeSymbols, Turing.MultiTapeTM, Turing.MultiTapeTM.step}
\difficulty{5}
Let $M$ be a Turing machine with $k$ work tapes, let $K$ satisfy $k\le K$, let $x$ be a
Boolean input, and let $c$ be a configuration of $M$ on $x$. One step of the machine
padded to $K$ work tapes, starting from the padded configuration, yields the padding of
the configuration obtained from $c$ by one step of $M$.
\end{lemma}

\begin{lemma}[Padding preserves machine runs]
\lean{Complexity.LogProg.padTM_runFrom}
\label{Complexity.LogProg.padTM_runFrom}
\leanok
\uses{Complexity.LogProg.padCfg, Complexity.LogProg.padTM, Complexity.LogProg.padTM_step, Turing.Cfg, Turing.MultiTapeTM, Turing.MultiTapeTM.runFrom, Turing.MultiTapeTM.runFrom_comm_of_step}
\difficulty{1}
Let $M$ be a Turing machine with $k$ work tapes, let $K$ be a natural number with $k\le
K$, let $x$ be a Boolean input, and let $c$ be a configuration of $M$ on $x$. For every
$n\in\mathbb{N}$, running the machine padded to $K$ work tapes for $n$ steps from the
padded configuration gives the padding of the configuration reached by running $M$ from
$c$ for $n$ steps.
\end{lemma}

\begin{definition}[Tagged machine family]
\lean{Complexity.LogProg.sigmaTM}
\label{Complexity.LogProg.sigmaTM}
\leanok
\uses{Turing.MultiTapeTM}
Let $d$ be a natural number, let $S_j$ be a state set for each $j
e d$ with $0\le j<d$, and let each $M_j$ be a deterministic machine with $K$ Boolean
work tapes and state set $S_j$. Define a machine with $K$ work tapes whose states are
either a distinguished initial state or a pair consisting of an index $j$ and a state in
$S_j$. The initial state is distinguished; from it the machine halts, while from a
tagged state it performs the corresponding machine's transition and tags any successor
state with the same index.
\end{definition}

\begin{definition}[Tagged configuration]
\lean{Complexity.LogProg.sigCfg}
\label{Complexity.LogProg.sigCfg}
\leanok
\uses{Turing.Cfg}
Let $d$ be a natural number, let $S_j$ be a state set for each $j<d$, and fix an input
string. Given an index $j<d$ and a configuration on $K$ Boolean work tapes with state
set $S_j$, define a configuration with tagged states: each present state $q$ becomes the
pair $(j,q)$, while a missing state remains missing. The input head position, work
tapes, work head positions, and output are unchanged.
\end{definition}

\begin{lemma}[One-step simulation by a tagged machine]
\lean{Complexity.LogProg.sigmaTM_step}
\label{Complexity.LogProg.sigmaTM_step}
\leanok
\uses{Complexity.LogProg.sigCfg, Complexity.LogProg.sigmaTM, Turing.Cfg, Turing.MultiTapeTM, Turing.MultiTapeTM.step}
\difficulty{4}
Let $d,K\in\bbn$, let $S_j$ be a state type for each of the $d$ indices, and let $M_j$
be a $K$-work-tape deterministic Turing machine with Boolean input and state type $S_j$.
For every Boolean input $x$, index $j<d$, and configuration $c$ of $M_j$ on $x$, one
step of the combined machine from the configuration tagging $c$ with $j$ equals the
configuration obtained by taking one step of $M_j$ from $c$ and tagging the result with
$j$.
\end{lemma}

\begin{lemma}[Simulation of a tagged machine run]
\lean{Complexity.LogProg.sigmaTM_runFrom}
\label{Complexity.LogProg.sigmaTM_runFrom}
\leanok
\uses{Complexity.LogProg.sigCfg, Complexity.LogProg.sigmaTM, Complexity.LogProg.sigmaTM_step, Turing.Cfg, Turing.MultiTapeTM, Turing.MultiTapeTM.runFrom, Turing.MultiTapeTM.runFrom_comm_of_step}
\difficulty{1}
Let $d,K\in\bbn$, let $S_j$ be a state type for each $j\in\{0,\ldots,d-1\}$, and let
$M_j$ be a $K$-work-tape deterministic Turing machine with Boolean input and state type
$S_j$. A tagged configuration places the state of $M_j$ inside the disjoint union of the
state types while leaving the tape contents, head positions, input-head position, and
output unchanged. For every Boolean input $x$, index $j$, configuration $c$ of $M_j$ on
$x$, and $n\in\bbn$, running the combined machine for $n$ steps from the tagged
configuration gives exactly the tagged result of running $M_j$ on $c$ for $n$ steps.
\end{lemma}

\begin{definition}[Maximum work tape count]
\lean{Complexity.LogProg.bankK}
\label{Complexity.LogProg.bankK}
\leanok
Given a natural number $d$ and a family of Boolean finite-tape machines indexed by
$j<d$, define the bank's work tape count to be the supremum of the work tape counts of
the machines in the family.
\end{definition}

\begin{lemma}[Work-Tape Bound for the Machine Bank]
\lean{Complexity.LogProg.le_bankK}
\label{Complexity.LogProg.le_bankK}
\leanok
\uses{Complexity.LogProg.bankK}
\difficulty{1}
Let $d$ be a natural number, and let a machine be assigned to each index in the finite
set $\text{Fin}(d)$. The number of work tapes of the machine at any such index is at most
the maximum number of work tapes among the machines in the collection.
\end{lemma}

\begin{definition}[Bank state set]
\lean{Complexity.LogProg.BankS}
\label{Complexity.LogProg.BankS}
\leanok
Given a natural number $d$ and a family of Boolean finite-tape machines indexed by
$j<d$, define the bank state type to consist of either no state or a pair $(j,q)$, where
$j<d$ and $q$ is a state of the $j$th machine.
\end{definition}

\begin{definition}[Bank simulation machine]
\lean{Complexity.LogProg.bankTM}
\label{Complexity.LogProg.bankTM}
\leanok
\uses{Complexity.LogProg.BankS, Complexity.LogProg.CleanSt, Complexity.LogProg.bankK, Complexity.LogProg.cleanTM, Complexity.LogProg.padTM, Complexity.LogProg.sigmaTM, Turing.MultiTapeTM}
Given a natural number $d$ and a family of Boolean finite-tape machines indexed by
$j<d$, define a machine with twice the bank's work tape count, Boolean alphabet, and
cleaning states associated with the bank state type and bank tape count. It is the
cleaning construction applied to the tagged machine family after padding each member to
the bank's work tape count, with the distinguished initial state as its starting state.
\end{definition}

\begin{definition}[Bank starting state]
\lean{Complexity.LogProg.bankStart}
\label{Complexity.LogProg.bankStart}
\leanok
\uses{Complexity.LogProg.BankS, Complexity.LogProg.CleanSt, Complexity.LogProg.bankK}
Given a natural number $d$, a family of Boolean finite-tape machines indexed by $j<d$,
and an index $j<d$, define the bank's starting cleaning state to be the initialization
state tagged with the pair consisting of $j$ and the initial state of the $j$th machine.
\end{definition}

\begin{theorem}[Space-bounded clean simulation of a machine bank]
\lean{Complexity.LogProg.bank_cleanRun}
\label{Complexity.LogProg.bank_cleanRun}
\leanok
\uses{Complexity.LogProg.BankS, Complexity.LogProg.CleanRun, Complexity.LogProg.bankK, Complexity.LogProg.bankStart, Complexity.LogProg.bankTM, Complexity.LogProg.cleanTM_run, Complexity.LogProg.le_bankK, Complexity.LogProg.padCfg, Complexity.LogProg.padTM, Complexity.LogProg.padTM_runFrom, Complexity.LogProg.sigCfg, Complexity.LogProg.sigmaTM, Complexity.LogProg.sigmaTM_runFrom, Turing.Cfg, Turing.Cfg.init, Turing.FinTM.DecidesInSpace, Turing.FinTM.computesInTime_iff, Turing.MultiTapeTM.indicator, Turing.MultiTapeTM.initCfg, Turing.MultiTapeTM.runFrom, Turing.MultiTapeTM.spaceUsedByTape_le_spaceUsed, Turing.MultiTapeTM.visitedByTapeHead}
\difficulty{6}
Let $d$ be a natural number and let $(M_j)_{j<d}$ be finite-state Boolean Turing
machines. For each $j<d$, let $A_j$ be a language of Boolean words and $s_j$ a space
bound, and suppose $M_j$ decides $A_j$ using at most $s_j(n)$ visited work-tape cells on
inputs of length $n$. Then, for every $j<d$ and Boolean input $V$, the bank machine
started in the state corresponding to $M_j$'s initial state halts with output consisting
of the single bit indicating whether $V\in A_j$, has all work tapes blank and all
work-tape heads at the origin, and throughout the run every head position has absolute
value at most $\max(s_j(|V|),1)$.
\end{theorem}
